\documentclass[11pt,a4paper]{article}
% =====================================================================
%  Gemeinsame Vorlage der Arbeitsblätter
%  "Objektorientierte Programmierung mit Java" – Informatik 10 (NTG)
%  Einbinden in jedem Blatt direkt nach \documentclass: % =====================================================================
%  Gemeinsame Vorlage der Arbeitsblätter
%  "Objektorientierte Programmierung mit Java" – Informatik 10 (NTG)
%  Einbinden in jedem Blatt direkt nach \documentclass: % =====================================================================
%  Gemeinsame Vorlage der Arbeitsblätter
%  "Objektorientierte Programmierung mit Java" – Informatik 10 (NTG)
%  Einbinden in jedem Blatt direkt nach \documentclass: \input{ab-vorlage}
%  Übersetzen mit XeLaTeX, LuaLaTeX, pdfLaTeX oder Tectonic.
%
%  Blatt:   tectonic ab-2-5-arrays.tex   (oder xelatex)
%  Lösung:  tectonic ab-2-5-arrays-lsg.tex
%           (Datei enthält nur: \def\mitloesung{}\input{ab-2-5-arrays};
%            füllt alle Lücken rot aus; siehe \lsg, \lsgrest, \lsglinien,
%            \lsgfeld, \lsgkariert, \lsgkarobox, \lsgtext, \lsgoder)
%  Lösungseinträge belegen nie mehr Platz als die Lücke im Blatt.
% =====================================================================
\usepackage[a4paper,left=18mm,right=18mm,top=26mm,bottom=17mm,headheight=15mm,headsep=4mm,footskip=8mm]{geometry}
\usepackage{iftex}
\ifPDFTeX
  \usepackage[T1]{fontenc}
  \usepackage[utf8]{inputenc}
  \usepackage{helvet}
  \usepackage[scaled=0.86]{beramono}
  \renewcommand{\familydefault}{\sfdefault}
\else
  \usepackage{fontspec}
  \setmainfont{texgyreheros}[Extension=.otf,UprightFont=*-regular,BoldFont=*-bold,ItalicFont=*-italic,BoldItalicFont=*-bolditalic]
  \setmonofont{DejaVuSansMono}[Extension=.ttf,UprightFont=*,BoldFont=*-Bold,ItalicFont=*-Oblique,BoldItalicFont=*-BoldOblique,Scale=0.86]
\fi
\usepackage[ngerman,shorthands=off]{babel}
\usepackage{xcolor,graphicx,tabularx,array,enumitem,fancyhdr,tikz,listings,multicol,calc,etoolbox,ragged2e,colortbl,xspace,pifont}
\usepackage[most]{tcolorbox}
\usetikzlibrary{arrows.meta,positioning,calc,decorations.pathreplacing}
\usepackage[hidelinks]{hyperref}
\setlength{\parindent}{0pt}
\setlength{\parskip}{3pt}
\setlist{nosep,leftmargin=*}

% ---------- Lösungsschalter ----------
\newif\ifloesung
\ifdefined\mitloesung\loesungtrue\fi
\newcommand{\lsgfarbe}{\color{red!75!black}}

% ---------- Kopf- und Fußzeile (einheitliche Form) ----------
\newcommand{\lehrkraft}{Hau}
\newcommand{\themenbereich}{2. Objektorientierte Programmierung}
\newcommand{\abnummer}{}
\pagestyle{fancy}\fancyhf{}
\renewcommand{\headrulewidth}{0.4pt}
\fancyhead[L]{\small Klasse 10\\Datum:}
\fancyhead[C]{\small Themenbereich:\\\themenbereich}
\fancyhead[R]{\small \lehrkraft\ifloesung\\{\lsgfarbe\bfseries Lösung}\fi}
\fancyfoot[C]{\scriptsize edu-mrh.de\,·\,Informatik 10\,·\,Blatt \abnummer\,·\,Seite \thepage}
\newcommand{\abtitel}[2]{\renewcommand{\abnummer}{#1}\begin{center}\Large\bfseries #1\quad\underline{#2}\end{center}\vspace{-1mm}}

% ---------- Boxen ----------
\newenvironment{aufgabe}[2][]{\begin{tcolorbox}[enhanced,breakable,colback=white,colframe=black,boxrule=0.7pt,arc=3mm,left=2mm,right=2mm,top=1.5mm,bottom=1.5mm,before skip=2.5mm,after skip=2.5mm]\textbf{Aufgabe #2}\ifstrempty{#1}{}{ \textit{#1}}\par\vspace{0.6mm}}{\end{tcolorbox}}
\newtcolorbox{merke}[1][Merke]{enhanced,breakable,colback=green!5,colframe=green!40!black,boxrule=0.8pt,arc=2mm,left=2mm,right=2mm,top=1.5mm,bottom=1.5mm,before skip=2.5mm,after skip=2.5mm,title={#1},fonttitle=\bfseries,coltitle=white}
\newtcolorbox{hinweis}{enhanced,breakable,colback=yellow!12,colframe=orange!70!black,boxrule=0.6pt,arc=2mm,left=2mm,right=2mm,top=1mm,bottom=1mm,before skip=2mm,after skip=2mm}
\newcommand{\web}[1]{{\small\color{gray}\textit{Website: #1}}}

% ---------- Lücken, Linien, Karos ----------
% Einheitliche Maße für alle Blätter der Sequenz:
%  \linien{n}       n Schreiblinien, Abstand 8 mm, erste Linie 8 mm unter dem Text
%  \kariert{n}      Karoraster 5 mm, n Kästchen hoch, nur ganze Kästchen, mittig
%  \karobox{b}{h}   Karoraster b x h, auf ganze Kästchen gerundet
%  \luecke[l]       Lücke im Fließtext (Standard 4,5 cm, Fachbegriffe mind. 3,5 cm)
%  (jeweils mit Lösungsvariante \lsg…, siehe Kopf der Datei)
%  \restzeile       Lücke bis zum Zeilenende
%  \lueckentext     am Anfang eines Absatzes/Kastens mit Lücken: größerer
%                   Zeilenabstand und Flattersatz, damit man in die Lücken schreiben kann
%  \linienfeld{n}   umrahmtes Feld mit n Schreiblinien (z. B. für Java-Code)
%  \schreibzeile    Strut für Tabellenzellen, in die geschrieben wird (8 mm hoch)
\newlength{\zeilenabstand}\setlength{\zeilenabstand}{8mm}
\newlength{\karo}\setlength{\karo}{5mm}
\tikzset{karo/.style={draw=gray!50,line width=0.3pt},schreiblinie/.style={draw=black!65,line width=0.35pt}}
\newlength{\lsgbreite}
% Lücke: \luecke[Breite] / mit Lösung \lsg[Breite]{Text} (zu langer Text wird verkleinert)
\newcommand{\lsg}[2][4.5cm]{\leavevmode\ifloesung\settowidth{\lsgbreite}{\,#2\,}%
  \ifdim\lsgbreite>#1\relax\rlap{\raisebox{0.8pt}{\resizebox{#1}{!}{\lsgfarbe\,#2\,}}}%
  \else\rlap{\raisebox{0.8pt}{\lsgfarbe\,#2}}\fi\fi\rule[-0.6ex]{#1}{0.4pt}\xspace}
\newcommand{\luecke}[1][4.5cm]{\lsg[#1]{}}
% Lücke bis Zeilenende: \restzeile / \lsgrest{Text}
\newcommand{\lsgrest}[1]{\leavevmode\unskip\hspace{1.5mm}\ifloesung\rlap{\raisebox{0.8pt}{\lsgfarbe\,#1}}\fi\leaders\hbox{\rule[-0.6ex]{0.5pt}{0.4pt}}\hfill\null\par}
\newcommand{\restzeile}{\lsgrest{}}
% nur in der Lösung sichtbar (Platz bleibt frei): \lsgtext{…}; Alternativen: \lsgoder{Lösung}{Blatt}
\newcommand{\lsgtext}[1]{\ifloesung\leavevmode{\lsgfarbe#1}\fi}
\newcommand{\lsgoder}[2]{\ifloesung\leavevmode#1\else#2\fi}
\newcommand{\lueckentext}{\linespread{1.5}\selectfont\raggedright}
\newcommand{\schreibzeile}{\rule[-2.8mm]{0pt}{8mm}}
% Schritt-Nummern wie im Code und auf der Website: \nr{1} = ①
\newcommand{\nr}[1]{\ding{\numexpr171+#1\relax}}
% Schreiblinien: \linien{n} / \lsglinien{n}{Text} (Text liegt auf den Linien)
\newcommand{\lsglinien}[2]{\par\vspace{0.5mm}\noindent\begin{tikzpicture}
  \path (0,0) (\linewidth-0.4pt,-#1\zeilenabstand-1.5mm);
  \foreach \i in {1,...,#1}{\draw[schreiblinie] (0.2pt,-\i\zeilenabstand) -- (\linewidth-0.2pt,-\i\zeilenabstand);}
  \ifloesung\node[overlay,anchor=north west,inner sep=0pt] at (1mm,-\zeilenabstand+1pt+5mm) {\parbox[t]{\linewidth-2mm}{\lsgfarbe\raggedright\setlength{\baselineskip}{\zeilenabstand}\rule{0pt}{5mm}#2}};\fi
  \end{tikzpicture}\par}
\newcommand{\linien}[1]{\lsglinien{#1}{}}
% Codefeld mit Linien: \linienfeld{n} / \lsgfeld{n}{Code, Zeilen mit \\ getrennt}
\newcommand{\lsgfeld}[2]{\noindent\begin{tikzpicture}
  \draw[line width=0.4pt] (0.2pt,0) rectangle (\linewidth-0.2pt,-#1\zeilenabstand-3mm);
  \foreach \i in {1,...,#1}{\draw[schreiblinie] (2mm,-\i\zeilenabstand) -- (\linewidth-2mm,-\i\zeilenabstand);}
  \ifloesung\node[overlay,anchor=north west,inner sep=0pt] at (2.5mm,-\zeilenabstand+1pt+5mm) {\parbox[t]{\linewidth-5mm}{\lsgfarbe\ttfamily\small\setlength{\baselineskip}{\zeilenabstand}\rule{0pt}{5mm}#2}};\fi
  \end{tikzpicture}}
\newcommand{\linienfeld}[1]{\lsgfeld{#1}{}}
% Karoraster; Lösungszeichnung als TikZ-Code, Ursprung oben links, y nach unten negativ:
% \kariert{n} / \lsgkariert{n}{TikZ},  \karobox{b}{h} / \lsgkarobox{b}{h}{TikZ}
\newcommand{\karoraster}[3]{\begin{tikzpicture}\draw[karo] (0,0) grid[step=\karo] (#1\karo,#2\karo);
  \ifloesung\begin{scope}[overlay,shift={(0,#2\karo)},red!75!black,text=red!75!black,line width=0.8pt]#3\end{scope}\fi\end{tikzpicture}}
\newcommand{\lsgkariert}[2]{\par\vspace{1.5mm}\noindent\pgfmathtruncatemacro{\karoX}{floor(\linewidth/\karo)}%
  \makebox[\linewidth]{\karoraster{\karoX}{#1}{#2}}\par\vspace{1mm}}
\newcommand{\kariert}[1]{\lsgkariert{#1}{}}
\newcommand{\lsgkarobox}[3]{\pgfmathtruncatemacro{\karoX}{round((#1)/\karo)}\pgfmathtruncatemacro{\karoY}{round((#2)/\karo)}\karoraster{\karoX}{\karoY}{#3}}
\newcommand{\karobox}[2]{\lsgkarobox{#1}{#2}{}}
% Lösungs-Klassenkarte (für Zeichnungen im Karo): \lsgkarte{Name}{Attribute}{Methoden}
\newcommand{\lsgkarte}[3]{{\lsgfarbe\setlength{\arrayrulewidth}{0.6pt}\arrayrulecolor{red!75!black}\begin{tabular}{|l|}\hline\textbf{#1}\\\hline #2\\\hline #3\\\hline\end{tabular}}}
% Lösungs-Objektkarte (in TikZ als Knoteninhalt): \node[objekt] {\lsgobjekt{name : Klasse}{attr = wert\\…}};
\newcommand{\lsgobjekt}[2]{{\lsgfarbe\small\arrayrulecolor{red!75!black}\begin{tabular}{@{}l@{}}\textbf{#1}\\\hline #2\end{tabular}}}
\tikzset{objekt/.style={draw,rounded corners=3mm,inner sep=2mm,anchor=north west}}
\newcommand{\klassenkarte}[3]{\begin{tabular}{|l|}\hline\textbf{#1}\\\hline #2\\\hline #3\\\hline\end{tabular}}
\newcommand{\code}[1]{\texttt{#1}}

% ---------- Java-Listings ----------
\lstset{language=Java,basicstyle=\ttfamily\small,keywordstyle=\bfseries,commentstyle=\color{gray!80!black},stringstyle=\color{black},showstringspaces=false,columns=flexible,keepspaces=true,tabsize=3,frame=single,framerule=0.4pt,rulecolor=\color{black},xleftmargin=2mm,framexleftmargin=1mm,aboveskip=2mm,belowskip=2mm,breaklines=true,
 literate={ä}{{\"a}}1 {ö}{{\"o}}1 {ü}{{\"u}}1 {Ä}{{\"A}}1 {Ö}{{\"O}}1 {Ü}{{\"U}}1 {ß}{{\ss}}1 {–}{{--}}1 {„}{{\glqq}}1 {“}{{\grqq}}1}
\lstnewenvironment{java}[1][]{\lstset{#1}}{}
% Lückencode: größerer Zeilenabstand, Lücken mit (*\luecke[3cm]*) im Code
\lstnewenvironment{javaluecke}[1][]{\lstset{basicstyle=\linespread{1.5}\small\ttfamily,escapeinside={(*}{*)},#1}}{}

%  Übersetzen mit XeLaTeX, LuaLaTeX, pdfLaTeX oder Tectonic.
%
%  Blatt:   tectonic ab-2-5-arrays.tex   (oder xelatex)
%  Lösung:  tectonic ab-2-5-arrays-lsg.tex
%           (Datei enthält nur: \def\mitloesung{}\documentclass[11pt,a4paper]{article}
\input{ab-vorlage}
\begin{document}
\abtitel{2.5}{Verwaltung gleichartiger Objekte -- Das Array / Feld}

\begin{aufgabe}[Projekt Zufallszahlen]{1}
\textbf{a)} Ergänze die Attribute \code{zahl1} bis \code{zahl3}.\quad \textbf{b)} Weise ihnen im Konstruktor Werte von 1 bis 6 zu.\\
\textbf{c)} Und bei 100 Zahlen? Skizziere, wie ein Speicher aussehen müsste, in dem du „die i-te Zahl“ ansprechen kannst:
\lsgkariert{4}{\begin{scope}[font=\scriptsize]
\foreach \v [count=\i from 0] in {5,2,6,1,4} {\draw (10mm+\i*10mm,-7mm) rectangle ++(10mm,-7mm); \node at (15mm+\i*10mm,-10.5mm) {\v}; \node at (15mm+\i*10mm,-4.5mm) {\i};}
\node at (65mm,-10.5mm) {\dots}; \draw (70mm,-7mm) rectangle ++(10mm,-7mm); \node at (75mm,-10.5mm) {3}; \node at (75mm,-4.5mm) {99};
\node[anchor=west] at (2mm,-4.5mm) {i:}; \node[anchor=west,align=left] at (85mm,-9mm) {ein Name für alle 100 Zahlen,\\Zugriff über die Nummer (Index)};
\end{scope}}
\end{aufgabe}

\begin{minipage}[t]{0.7\linewidth}
\begin{merke}[Felder]\lueckentext
Felder (engl. \lsg[3cm]{Arrays}) dienen der Speicherung einer \lsg[3cm]{festen} \lsg[3cm]{Anzahl} von Werten \lsg[4.5cm]{gleichen Datentyps}, die über einen Index (Positionszahl) adressiert werden können.

Dadurch kann z.\,B. mit einer Wiederholung auf alle Elemente zugegriffen werden, ohne jedes einzeln zu behandeln.
\end{merke}
\end{minipage}\hfill
\begin{minipage}[t]{0.27\linewidth}
\vspace{3mm}
\begin{tabular}{|c|c|}\hline
Index i & \code{textFeld[i]}\\\hline
0 & \code{"Hallo "}\\\hline
1 & \code{"Welt, "}\\\hline
2 & \code{"ein "}\\\hline
3 & \code{"Satz."}\\\hline
\end{tabular}
\end{minipage}

\begin{aufgabe}{2}
Übertrage den Code in das Projekt (TODO 2) und beschreibe die Funktion der Zeilen.
\end{aufgabe}
\begin{javaluecke}
class Zufallszahlen {
   private int[] zahlenFeld;              // (*\lsg[7.5cm]{Deklaration: Feld für int-Werte}*)
   Zufallszahlen() {
      zahlenFeld = new int[8];            // (*\lsg[7.5cm]{Erzeugung: 8 Zellen, alle mit 0 belegt}*)
      zahlenFeld[0] = 5;                  // (*\lsg[7.5cm]{erste Zelle (Index 0) bekommt 5}*)
      zahlenFeld[7] = Random.randint(1, 6); // (*\lsg[7.1cm]{letzte Zelle bekommt Zufallszahl 1--6}*)
   }
}
\end{javaluecke}

\begin{minipage}[t]{0.7\linewidth}
\begin{merke}\lueckentext
\begin{itemize}
\item Felder beginnen ihren Index immer bei \lsg[1.5cm]{0}.
\item Die Größe eines Feldes liefert \code{<Feldname>.length}. Hier also: \lsg[4cm]{\code{zahlenFeld.length}}. Das Array ist \lsg[1.5cm]{8} Zellen groß.
\item Was passiert bei \code{zahlenFeld[8]}?\lsgrest{Laufzeitfehler (Index zu groß)}
\end{itemize}
\end{merke}
\end{minipage}\hfill
\begin{minipage}[t]{0.27\linewidth}
\vspace{3mm}
\begin{tabular}{|c|c|}\hline
Index i & \code{zahlenFeld[i]}\\\hline
0 & 5\\\hline
1 & 0\\\hline
\rule[-1.2mm]{0pt}{5.2mm}2 & \lsgtext{0}\\\hline
\rule[-1.2mm]{0pt}{5.2mm}\lsgtext{3} & \lsgtext{0}\\\hline
\rule[-1.2mm]{0pt}{5.2mm}\lsgtext{4} & \lsgtext{0}\\\hline
\rule[-1.2mm]{0pt}{5.2mm}\lsgtext{5} & \lsgtext{0}\\\hline
\rule[-1.2mm]{0pt}{5.2mm}\lsgtext{6} & \lsgtext{0}\\\hline
\rule[-1.2mm]{0pt}{5.2mm}\lsgtext{7} & \lsgtext{1 bis 6}\\\hline
\end{tabular}
\end{minipage}

\newpage
Fülle die Tabelle auf Seite 1. Warum ist \code{zahlenFeld[7]} die letzte Zelle?
\lsglinien{2}{Der Index beginnt bei 0. Bei 8 Zellen laufen die Indizes von 0 bis 7, der letzte ist also 8 $-$ 1 = 7.}
{\lueckentext Letzte Zelle bei beliebiger Feldgröße:\lsgrest{\code{zahlenFeld[zahlenFeld.length - 1]}}}

\begin{aufgabe}{3}
\textbf{a)} Weise allen Zellen zufällige Werte zu (Wiederholung mit fester Anzahl).\quad
\textbf{b)} 100 Zellen, jede mit der Summe zweier Würfel.\\
Was änderst du an der Schleife aus a), damit sie bei jeder Feldgröße stimmt?
\lsglinien{1}{Bedingung \code{i < zahlenFeld.length} statt einer festen Zahl.}
\end{aufgabe}

\textbf{1:n-Beziehung als Sammlung umsetzen -- Arrays mit Referenzen}

{\centering
\begin{tikzpicture}[x=1cm,y=1cm]
\node[draw,fill=cyan!15,minimum height=6mm] (m) at (0,0) {Mannschaft};
\node[draw,fill=cyan!15,minimum height=6mm] (s) at (5,0) {Spieler};
\draw (m) -- (s) node[midway,above] {\small $<$ spielt in} node[pos=0.08,below] {\small 1} node[pos=0.92,below] {\small n};
\foreach \i/\n in {0/paul,1/tim,2/maja,3/tina} {
  \node[draw,fill=red!10,minimum width=13mm,minimum height=5mm] (z\i) at (6.6+\i*2.2,0.3) {\small\bfseries \i};
  \node[draw,fill=red!10,font=\scriptsize\ttfamily,inner sep=2pt] (o\i) at (6.6+\i*2.2,-1.1) {\n: Spieler};
  \draw[-{Stealth}] (z\i.south) -- (o\i.north);
}
\end{tikzpicture}
\par}
{\lueckentext Auch Referenzen können in Arrays verwaltet werden. Dabei verweist jede Zelle auf ein \lsg[3.5cm]{Objekt} der anderen \lsg[3.5cm]{Klasse}, hier: \lsg[3.5cm]{\code{Spieler}}.\par}

\begin{hinweis}
\lueckentext\textbf{Beachte:} Alle Elemente eines Feldes müssen denselben \lsg[3.5cm]{Datentyp} besitzen, dieser kann aber beliebig gewählt werden: auch \code{boolean}, \code{String}, \code{Rectangle}, \ldots
\end{hinweis}

\begin{merke}[Syntaxübersicht]\small
\begin{tabular}{@{}l l@{}}
\code{<Sichtbarkeit> <Datentyp>[] <Feldname>;} & Deklaration: \code{[]} nach dem Datentyp macht ein Feld daraus.\\
\code{<Feldname> = new <Datentyp>[<Groesse>];} & Erzeugung mit fester Größe (nicht veränderbar).\\
\code{<Feldname>[<Index>]} & Zugriff auf die Zelle an dieser Position.\\
\end{tabular}

\textbf{Beispiele:}\\
lesen: \code{int eineZahl = zahlenFeld[0];}\hfill schreiben: \code{zahlenFeld[0] = eineZahl;}\\
Methode aufrufen: \code{kreisFeld[0].tuEtwas();}\hfill Attribut: \code{kreisFeld[0].xPosition = 100;}
\end{merke}

\begin{minipage}[t]{0.74\linewidth}
\begin{merke}[Mehrdimensionale Arrays]\lueckentext
Arrays können mehrere Dimensionen haben und so eine \lsg[3.5cm]{Tabelle (Raster)} aufspannen.\\
Deklaration und Erzeugung:\lsgrest{\code{Rectangle[][] fuellung}}
\lsgrest{\code{fuellung = new Rectangle[5][5];}}
Zugriff auf Spalte x, Zeile y:\lsgrest{\code{fuellung[x][y]}}
\end{merke}
\end{minipage}\hfill
\begin{minipage}[t]{0.23\linewidth}
\vspace{1mm}\centering
\begin{tikzpicture}[x=6mm,y=6mm,font=\scriptsize]
\ifloesung\fill[red!25] \foreach \x/\y in {1/1,3/1,1/3,2/4,3/3} {(\x,-\y) rectangle (\x+1,-\y-1)};\fi
\draw[step=6mm] (0,0) grid (5,-5);
\foreach \i in {0,...,4} {\node at (\i+0.5,0.4) {\i}; \node at (-0.4,-\i-0.5) {\i};}
\node at (2.5,1.05) {\code{x} $\rightarrow$}; \node[rotate=-90] at (-1.05,-2.5) {\code{y} $\rightarrow$};
\end{tikzpicture}
\end{minipage}

\begin{aufgabe}[Projekt Mosaik]{4}
\textbf{a)} TODO 1 und 2 in \code{Mosaik1}: zweidimensionales Array passender Größe erzeugen und in \code{neuesMosaik()} füllen.\\
\textbf{b)} Auf \code{Mosaik2} wechseln, TODO 3--5 (Smiley). Markiere die Smiley-Zellen \code{fuellung[x][y]} im Raster oben.\lsgtext{\ (z.\,B. [1][1], [3][1], [1][3], [2][4], [3][3])}
\end{aufgabe}
\end{document}
;
%            füllt alle Lücken rot aus; siehe \lsg, \lsgrest, \lsglinien,
%            \lsgfeld, \lsgkariert, \lsgkarobox, \lsgtext, \lsgoder)
%  Lösungseinträge belegen nie mehr Platz als die Lücke im Blatt.
% =====================================================================
\usepackage[a4paper,left=18mm,right=18mm,top=26mm,bottom=17mm,headheight=15mm,headsep=4mm,footskip=8mm]{geometry}
\usepackage{iftex}
\ifPDFTeX
  \usepackage[T1]{fontenc}
  \usepackage[utf8]{inputenc}
  \usepackage{helvet}
  \usepackage[scaled=0.86]{beramono}
  \renewcommand{\familydefault}{\sfdefault}
\else
  \usepackage{fontspec}
  \setmainfont{texgyreheros}[Extension=.otf,UprightFont=*-regular,BoldFont=*-bold,ItalicFont=*-italic,BoldItalicFont=*-bolditalic]
  \setmonofont{DejaVuSansMono}[Extension=.ttf,UprightFont=*,BoldFont=*-Bold,ItalicFont=*-Oblique,BoldItalicFont=*-BoldOblique,Scale=0.86]
\fi
\usepackage[ngerman,shorthands=off]{babel}
\usepackage{xcolor,graphicx,tabularx,array,enumitem,fancyhdr,tikz,listings,multicol,calc,etoolbox,ragged2e,colortbl,xspace,pifont}
\usepackage[most]{tcolorbox}
\usetikzlibrary{arrows.meta,positioning,calc,decorations.pathreplacing}
\usepackage[hidelinks]{hyperref}
\setlength{\parindent}{0pt}
\setlength{\parskip}{3pt}
\setlist{nosep,leftmargin=*}

% ---------- Lösungsschalter ----------
\newif\ifloesung
\ifdefined\mitloesung\loesungtrue\fi
\newcommand{\lsgfarbe}{\color{red!75!black}}

% ---------- Kopf- und Fußzeile (einheitliche Form) ----------
\newcommand{\lehrkraft}{Hau}
\newcommand{\themenbereich}{2. Objektorientierte Programmierung}
\newcommand{\abnummer}{}
\pagestyle{fancy}\fancyhf{}
\renewcommand{\headrulewidth}{0.4pt}
\fancyhead[L]{\small Klasse 10\\Datum:}
\fancyhead[C]{\small Themenbereich:\\\themenbereich}
\fancyhead[R]{\small \lehrkraft\ifloesung\\{\lsgfarbe\bfseries Lösung}\fi}
\fancyfoot[C]{\scriptsize edu-mrh.de\,·\,Informatik 10\,·\,Blatt \abnummer\,·\,Seite \thepage}
\newcommand{\abtitel}[2]{\renewcommand{\abnummer}{#1}\begin{center}\Large\bfseries #1\quad\underline{#2}\end{center}\vspace{-1mm}}

% ---------- Boxen ----------
\newenvironment{aufgabe}[2][]{\begin{tcolorbox}[enhanced,breakable,colback=white,colframe=black,boxrule=0.7pt,arc=3mm,left=2mm,right=2mm,top=1.5mm,bottom=1.5mm,before skip=2.5mm,after skip=2.5mm]\textbf{Aufgabe #2}\ifstrempty{#1}{}{ \textit{#1}}\par\vspace{0.6mm}}{\end{tcolorbox}}
\newtcolorbox{merke}[1][Merke]{enhanced,breakable,colback=green!5,colframe=green!40!black,boxrule=0.8pt,arc=2mm,left=2mm,right=2mm,top=1.5mm,bottom=1.5mm,before skip=2.5mm,after skip=2.5mm,title={#1},fonttitle=\bfseries,coltitle=white}
\newtcolorbox{hinweis}{enhanced,breakable,colback=yellow!12,colframe=orange!70!black,boxrule=0.6pt,arc=2mm,left=2mm,right=2mm,top=1mm,bottom=1mm,before skip=2mm,after skip=2mm}
\newcommand{\web}[1]{{\small\color{gray}\textit{Website: #1}}}

% ---------- Lücken, Linien, Karos ----------
% Einheitliche Maße für alle Blätter der Sequenz:
%  \linien{n}       n Schreiblinien, Abstand 8 mm, erste Linie 8 mm unter dem Text
%  \kariert{n}      Karoraster 5 mm, n Kästchen hoch, nur ganze Kästchen, mittig
%  \karobox{b}{h}   Karoraster b x h, auf ganze Kästchen gerundet
%  \luecke[l]       Lücke im Fließtext (Standard 4,5 cm, Fachbegriffe mind. 3,5 cm)
%  (jeweils mit Lösungsvariante \lsg…, siehe Kopf der Datei)
%  \restzeile       Lücke bis zum Zeilenende
%  \lueckentext     am Anfang eines Absatzes/Kastens mit Lücken: größerer
%                   Zeilenabstand und Flattersatz, damit man in die Lücken schreiben kann
%  \linienfeld{n}   umrahmtes Feld mit n Schreiblinien (z. B. für Java-Code)
%  \schreibzeile    Strut für Tabellenzellen, in die geschrieben wird (8 mm hoch)
\newlength{\zeilenabstand}\setlength{\zeilenabstand}{8mm}
\newlength{\karo}\setlength{\karo}{5mm}
\tikzset{karo/.style={draw=gray!50,line width=0.3pt},schreiblinie/.style={draw=black!65,line width=0.35pt}}
\newlength{\lsgbreite}
% Lücke: \luecke[Breite] / mit Lösung \lsg[Breite]{Text} (zu langer Text wird verkleinert)
\newcommand{\lsg}[2][4.5cm]{\leavevmode\ifloesung\settowidth{\lsgbreite}{\,#2\,}%
  \ifdim\lsgbreite>#1\relax\rlap{\raisebox{0.8pt}{\resizebox{#1}{!}{\lsgfarbe\,#2\,}}}%
  \else\rlap{\raisebox{0.8pt}{\lsgfarbe\,#2}}\fi\fi\rule[-0.6ex]{#1}{0.4pt}\xspace}
\newcommand{\luecke}[1][4.5cm]{\lsg[#1]{}}
% Lücke bis Zeilenende: \restzeile / \lsgrest{Text}
\newcommand{\lsgrest}[1]{\leavevmode\unskip\hspace{1.5mm}\ifloesung\rlap{\raisebox{0.8pt}{\lsgfarbe\,#1}}\fi\leaders\hbox{\rule[-0.6ex]{0.5pt}{0.4pt}}\hfill\null\par}
\newcommand{\restzeile}{\lsgrest{}}
% nur in der Lösung sichtbar (Platz bleibt frei): \lsgtext{…}; Alternativen: \lsgoder{Lösung}{Blatt}
\newcommand{\lsgtext}[1]{\ifloesung\leavevmode{\lsgfarbe#1}\fi}
\newcommand{\lsgoder}[2]{\ifloesung\leavevmode#1\else#2\fi}
\newcommand{\lueckentext}{\linespread{1.5}\selectfont\raggedright}
\newcommand{\schreibzeile}{\rule[-2.8mm]{0pt}{8mm}}
% Schritt-Nummern wie im Code und auf der Website: \nr{1} = ①
\newcommand{\nr}[1]{\ding{\numexpr171+#1\relax}}
% Schreiblinien: \linien{n} / \lsglinien{n}{Text} (Text liegt auf den Linien)
\newcommand{\lsglinien}[2]{\par\vspace{0.5mm}\noindent\begin{tikzpicture}
  \path (0,0) (\linewidth-0.4pt,-#1\zeilenabstand-1.5mm);
  \foreach \i in {1,...,#1}{\draw[schreiblinie] (0.2pt,-\i\zeilenabstand) -- (\linewidth-0.2pt,-\i\zeilenabstand);}
  \ifloesung\node[overlay,anchor=north west,inner sep=0pt] at (1mm,-\zeilenabstand+1pt+5mm) {\parbox[t]{\linewidth-2mm}{\lsgfarbe\raggedright\setlength{\baselineskip}{\zeilenabstand}\rule{0pt}{5mm}#2}};\fi
  \end{tikzpicture}\par}
\newcommand{\linien}[1]{\lsglinien{#1}{}}
% Codefeld mit Linien: \linienfeld{n} / \lsgfeld{n}{Code, Zeilen mit \\ getrennt}
\newcommand{\lsgfeld}[2]{\noindent\begin{tikzpicture}
  \draw[line width=0.4pt] (0.2pt,0) rectangle (\linewidth-0.2pt,-#1\zeilenabstand-3mm);
  \foreach \i in {1,...,#1}{\draw[schreiblinie] (2mm,-\i\zeilenabstand) -- (\linewidth-2mm,-\i\zeilenabstand);}
  \ifloesung\node[overlay,anchor=north west,inner sep=0pt] at (2.5mm,-\zeilenabstand+1pt+5mm) {\parbox[t]{\linewidth-5mm}{\lsgfarbe\ttfamily\small\setlength{\baselineskip}{\zeilenabstand}\rule{0pt}{5mm}#2}};\fi
  \end{tikzpicture}}
\newcommand{\linienfeld}[1]{\lsgfeld{#1}{}}
% Karoraster; Lösungszeichnung als TikZ-Code, Ursprung oben links, y nach unten negativ:
% \kariert{n} / \lsgkariert{n}{TikZ},  \karobox{b}{h} / \lsgkarobox{b}{h}{TikZ}
\newcommand{\karoraster}[3]{\begin{tikzpicture}\draw[karo] (0,0) grid[step=\karo] (#1\karo,#2\karo);
  \ifloesung\begin{scope}[overlay,shift={(0,#2\karo)},red!75!black,text=red!75!black,line width=0.8pt]#3\end{scope}\fi\end{tikzpicture}}
\newcommand{\lsgkariert}[2]{\par\vspace{1.5mm}\noindent\pgfmathtruncatemacro{\karoX}{floor(\linewidth/\karo)}%
  \makebox[\linewidth]{\karoraster{\karoX}{#1}{#2}}\par\vspace{1mm}}
\newcommand{\kariert}[1]{\lsgkariert{#1}{}}
\newcommand{\lsgkarobox}[3]{\pgfmathtruncatemacro{\karoX}{round((#1)/\karo)}\pgfmathtruncatemacro{\karoY}{round((#2)/\karo)}\karoraster{\karoX}{\karoY}{#3}}
\newcommand{\karobox}[2]{\lsgkarobox{#1}{#2}{}}
% Lösungs-Klassenkarte (für Zeichnungen im Karo): \lsgkarte{Name}{Attribute}{Methoden}
\newcommand{\lsgkarte}[3]{{\lsgfarbe\setlength{\arrayrulewidth}{0.6pt}\arrayrulecolor{red!75!black}\begin{tabular}{|l|}\hline\textbf{#1}\\\hline #2\\\hline #3\\\hline\end{tabular}}}
% Lösungs-Objektkarte (in TikZ als Knoteninhalt): \node[objekt] {\lsgobjekt{name : Klasse}{attr = wert\\…}};
\newcommand{\lsgobjekt}[2]{{\lsgfarbe\small\arrayrulecolor{red!75!black}\begin{tabular}{@{}l@{}}\textbf{#1}\\\hline #2\end{tabular}}}
\tikzset{objekt/.style={draw,rounded corners=3mm,inner sep=2mm,anchor=north west}}
\newcommand{\klassenkarte}[3]{\begin{tabular}{|l|}\hline\textbf{#1}\\\hline #2\\\hline #3\\\hline\end{tabular}}
\newcommand{\code}[1]{\texttt{#1}}

% ---------- Java-Listings ----------
\lstset{language=Java,basicstyle=\ttfamily\small,keywordstyle=\bfseries,commentstyle=\color{gray!80!black},stringstyle=\color{black},showstringspaces=false,columns=flexible,keepspaces=true,tabsize=3,frame=single,framerule=0.4pt,rulecolor=\color{black},xleftmargin=2mm,framexleftmargin=1mm,aboveskip=2mm,belowskip=2mm,breaklines=true,
 literate={ä}{{\"a}}1 {ö}{{\"o}}1 {ü}{{\"u}}1 {Ä}{{\"A}}1 {Ö}{{\"O}}1 {Ü}{{\"U}}1 {ß}{{\ss}}1 {–}{{--}}1 {„}{{\glqq}}1 {“}{{\grqq}}1}
\lstnewenvironment{java}[1][]{\lstset{#1}}{}
% Lückencode: größerer Zeilenabstand, Lücken mit (*\luecke[3cm]*) im Code
\lstnewenvironment{javaluecke}[1][]{\lstset{basicstyle=\linespread{1.5}\small\ttfamily,escapeinside={(*}{*)},#1}}{}

%  Übersetzen mit XeLaTeX, LuaLaTeX, pdfLaTeX oder Tectonic.
%
%  Blatt:   tectonic ab-2-5-arrays.tex   (oder xelatex)
%  Lösung:  tectonic ab-2-5-arrays-lsg.tex
%           (Datei enthält nur: \def\mitloesung{}\documentclass[11pt,a4paper]{article}
% =====================================================================
%  Gemeinsame Vorlage der Arbeitsblätter
%  "Objektorientierte Programmierung mit Java" – Informatik 10 (NTG)
%  Einbinden in jedem Blatt direkt nach \documentclass: \input{ab-vorlage}
%  Übersetzen mit XeLaTeX, LuaLaTeX, pdfLaTeX oder Tectonic.
%
%  Blatt:   tectonic ab-2-5-arrays.tex   (oder xelatex)
%  Lösung:  tectonic ab-2-5-arrays-lsg.tex
%           (Datei enthält nur: \def\mitloesung{}\input{ab-2-5-arrays};
%            füllt alle Lücken rot aus; siehe \lsg, \lsgrest, \lsglinien,
%            \lsgfeld, \lsgkariert, \lsgkarobox, \lsgtext, \lsgoder)
%  Lösungseinträge belegen nie mehr Platz als die Lücke im Blatt.
% =====================================================================
\usepackage[a4paper,left=18mm,right=18mm,top=26mm,bottom=17mm,headheight=15mm,headsep=4mm,footskip=8mm]{geometry}
\usepackage{iftex}
\ifPDFTeX
  \usepackage[T1]{fontenc}
  \usepackage[utf8]{inputenc}
  \usepackage{helvet}
  \usepackage[scaled=0.86]{beramono}
  \renewcommand{\familydefault}{\sfdefault}
\else
  \usepackage{fontspec}
  \setmainfont{texgyreheros}[Extension=.otf,UprightFont=*-regular,BoldFont=*-bold,ItalicFont=*-italic,BoldItalicFont=*-bolditalic]
  \setmonofont{DejaVuSansMono}[Extension=.ttf,UprightFont=*,BoldFont=*-Bold,ItalicFont=*-Oblique,BoldItalicFont=*-BoldOblique,Scale=0.86]
\fi
\usepackage[ngerman,shorthands=off]{babel}
\usepackage{xcolor,graphicx,tabularx,array,enumitem,fancyhdr,tikz,listings,multicol,calc,etoolbox,ragged2e,colortbl,xspace,pifont}
\usepackage[most]{tcolorbox}
\usetikzlibrary{arrows.meta,positioning,calc,decorations.pathreplacing}
\usepackage[hidelinks]{hyperref}
\setlength{\parindent}{0pt}
\setlength{\parskip}{3pt}
\setlist{nosep,leftmargin=*}

% ---------- Lösungsschalter ----------
\newif\ifloesung
\ifdefined\mitloesung\loesungtrue\fi
\newcommand{\lsgfarbe}{\color{red!75!black}}

% ---------- Kopf- und Fußzeile (einheitliche Form) ----------
\newcommand{\lehrkraft}{Hau}
\newcommand{\themenbereich}{2. Objektorientierte Programmierung}
\newcommand{\abnummer}{}
\pagestyle{fancy}\fancyhf{}
\renewcommand{\headrulewidth}{0.4pt}
\fancyhead[L]{\small Klasse 10\\Datum:}
\fancyhead[C]{\small Themenbereich:\\\themenbereich}
\fancyhead[R]{\small \lehrkraft\ifloesung\\{\lsgfarbe\bfseries Lösung}\fi}
\fancyfoot[C]{\scriptsize edu-mrh.de\,·\,Informatik 10\,·\,Blatt \abnummer\,·\,Seite \thepage}
\newcommand{\abtitel}[2]{\renewcommand{\abnummer}{#1}\begin{center}\Large\bfseries #1\quad\underline{#2}\end{center}\vspace{-1mm}}

% ---------- Boxen ----------
\newenvironment{aufgabe}[2][]{\begin{tcolorbox}[enhanced,breakable,colback=white,colframe=black,boxrule=0.7pt,arc=3mm,left=2mm,right=2mm,top=1.5mm,bottom=1.5mm,before skip=2.5mm,after skip=2.5mm]\textbf{Aufgabe #2}\ifstrempty{#1}{}{ \textit{#1}}\par\vspace{0.6mm}}{\end{tcolorbox}}
\newtcolorbox{merke}[1][Merke]{enhanced,breakable,colback=green!5,colframe=green!40!black,boxrule=0.8pt,arc=2mm,left=2mm,right=2mm,top=1.5mm,bottom=1.5mm,before skip=2.5mm,after skip=2.5mm,title={#1},fonttitle=\bfseries,coltitle=white}
\newtcolorbox{hinweis}{enhanced,breakable,colback=yellow!12,colframe=orange!70!black,boxrule=0.6pt,arc=2mm,left=2mm,right=2mm,top=1mm,bottom=1mm,before skip=2mm,after skip=2mm}
\newcommand{\web}[1]{{\small\color{gray}\textit{Website: #1}}}

% ---------- Lücken, Linien, Karos ----------
% Einheitliche Maße für alle Blätter der Sequenz:
%  \linien{n}       n Schreiblinien, Abstand 8 mm, erste Linie 8 mm unter dem Text
%  \kariert{n}      Karoraster 5 mm, n Kästchen hoch, nur ganze Kästchen, mittig
%  \karobox{b}{h}   Karoraster b x h, auf ganze Kästchen gerundet
%  \luecke[l]       Lücke im Fließtext (Standard 4,5 cm, Fachbegriffe mind. 3,5 cm)
%  (jeweils mit Lösungsvariante \lsg…, siehe Kopf der Datei)
%  \restzeile       Lücke bis zum Zeilenende
%  \lueckentext     am Anfang eines Absatzes/Kastens mit Lücken: größerer
%                   Zeilenabstand und Flattersatz, damit man in die Lücken schreiben kann
%  \linienfeld{n}   umrahmtes Feld mit n Schreiblinien (z. B. für Java-Code)
%  \schreibzeile    Strut für Tabellenzellen, in die geschrieben wird (8 mm hoch)
\newlength{\zeilenabstand}\setlength{\zeilenabstand}{8mm}
\newlength{\karo}\setlength{\karo}{5mm}
\tikzset{karo/.style={draw=gray!50,line width=0.3pt},schreiblinie/.style={draw=black!65,line width=0.35pt}}
\newlength{\lsgbreite}
% Lücke: \luecke[Breite] / mit Lösung \lsg[Breite]{Text} (zu langer Text wird verkleinert)
\newcommand{\lsg}[2][4.5cm]{\leavevmode\ifloesung\settowidth{\lsgbreite}{\,#2\,}%
  \ifdim\lsgbreite>#1\relax\rlap{\raisebox{0.8pt}{\resizebox{#1}{!}{\lsgfarbe\,#2\,}}}%
  \else\rlap{\raisebox{0.8pt}{\lsgfarbe\,#2}}\fi\fi\rule[-0.6ex]{#1}{0.4pt}\xspace}
\newcommand{\luecke}[1][4.5cm]{\lsg[#1]{}}
% Lücke bis Zeilenende: \restzeile / \lsgrest{Text}
\newcommand{\lsgrest}[1]{\leavevmode\unskip\hspace{1.5mm}\ifloesung\rlap{\raisebox{0.8pt}{\lsgfarbe\,#1}}\fi\leaders\hbox{\rule[-0.6ex]{0.5pt}{0.4pt}}\hfill\null\par}
\newcommand{\restzeile}{\lsgrest{}}
% nur in der Lösung sichtbar (Platz bleibt frei): \lsgtext{…}; Alternativen: \lsgoder{Lösung}{Blatt}
\newcommand{\lsgtext}[1]{\ifloesung\leavevmode{\lsgfarbe#1}\fi}
\newcommand{\lsgoder}[2]{\ifloesung\leavevmode#1\else#2\fi}
\newcommand{\lueckentext}{\linespread{1.5}\selectfont\raggedright}
\newcommand{\schreibzeile}{\rule[-2.8mm]{0pt}{8mm}}
% Schritt-Nummern wie im Code und auf der Website: \nr{1} = ①
\newcommand{\nr}[1]{\ding{\numexpr171+#1\relax}}
% Schreiblinien: \linien{n} / \lsglinien{n}{Text} (Text liegt auf den Linien)
\newcommand{\lsglinien}[2]{\par\vspace{0.5mm}\noindent\begin{tikzpicture}
  \path (0,0) (\linewidth-0.4pt,-#1\zeilenabstand-1.5mm);
  \foreach \i in {1,...,#1}{\draw[schreiblinie] (0.2pt,-\i\zeilenabstand) -- (\linewidth-0.2pt,-\i\zeilenabstand);}
  \ifloesung\node[overlay,anchor=north west,inner sep=0pt] at (1mm,-\zeilenabstand+1pt+5mm) {\parbox[t]{\linewidth-2mm}{\lsgfarbe\raggedright\setlength{\baselineskip}{\zeilenabstand}\rule{0pt}{5mm}#2}};\fi
  \end{tikzpicture}\par}
\newcommand{\linien}[1]{\lsglinien{#1}{}}
% Codefeld mit Linien: \linienfeld{n} / \lsgfeld{n}{Code, Zeilen mit \\ getrennt}
\newcommand{\lsgfeld}[2]{\noindent\begin{tikzpicture}
  \draw[line width=0.4pt] (0.2pt,0) rectangle (\linewidth-0.2pt,-#1\zeilenabstand-3mm);
  \foreach \i in {1,...,#1}{\draw[schreiblinie] (2mm,-\i\zeilenabstand) -- (\linewidth-2mm,-\i\zeilenabstand);}
  \ifloesung\node[overlay,anchor=north west,inner sep=0pt] at (2.5mm,-\zeilenabstand+1pt+5mm) {\parbox[t]{\linewidth-5mm}{\lsgfarbe\ttfamily\small\setlength{\baselineskip}{\zeilenabstand}\rule{0pt}{5mm}#2}};\fi
  \end{tikzpicture}}
\newcommand{\linienfeld}[1]{\lsgfeld{#1}{}}
% Karoraster; Lösungszeichnung als TikZ-Code, Ursprung oben links, y nach unten negativ:
% \kariert{n} / \lsgkariert{n}{TikZ},  \karobox{b}{h} / \lsgkarobox{b}{h}{TikZ}
\newcommand{\karoraster}[3]{\begin{tikzpicture}\draw[karo] (0,0) grid[step=\karo] (#1\karo,#2\karo);
  \ifloesung\begin{scope}[overlay,shift={(0,#2\karo)},red!75!black,text=red!75!black,line width=0.8pt]#3\end{scope}\fi\end{tikzpicture}}
\newcommand{\lsgkariert}[2]{\par\vspace{1.5mm}\noindent\pgfmathtruncatemacro{\karoX}{floor(\linewidth/\karo)}%
  \makebox[\linewidth]{\karoraster{\karoX}{#1}{#2}}\par\vspace{1mm}}
\newcommand{\kariert}[1]{\lsgkariert{#1}{}}
\newcommand{\lsgkarobox}[3]{\pgfmathtruncatemacro{\karoX}{round((#1)/\karo)}\pgfmathtruncatemacro{\karoY}{round((#2)/\karo)}\karoraster{\karoX}{\karoY}{#3}}
\newcommand{\karobox}[2]{\lsgkarobox{#1}{#2}{}}
% Lösungs-Klassenkarte (für Zeichnungen im Karo): \lsgkarte{Name}{Attribute}{Methoden}
\newcommand{\lsgkarte}[3]{{\lsgfarbe\setlength{\arrayrulewidth}{0.6pt}\arrayrulecolor{red!75!black}\begin{tabular}{|l|}\hline\textbf{#1}\\\hline #2\\\hline #3\\\hline\end{tabular}}}
% Lösungs-Objektkarte (in TikZ als Knoteninhalt): \node[objekt] {\lsgobjekt{name : Klasse}{attr = wert\\…}};
\newcommand{\lsgobjekt}[2]{{\lsgfarbe\small\arrayrulecolor{red!75!black}\begin{tabular}{@{}l@{}}\textbf{#1}\\\hline #2\end{tabular}}}
\tikzset{objekt/.style={draw,rounded corners=3mm,inner sep=2mm,anchor=north west}}
\newcommand{\klassenkarte}[3]{\begin{tabular}{|l|}\hline\textbf{#1}\\\hline #2\\\hline #3\\\hline\end{tabular}}
\newcommand{\code}[1]{\texttt{#1}}

% ---------- Java-Listings ----------
\lstset{language=Java,basicstyle=\ttfamily\small,keywordstyle=\bfseries,commentstyle=\color{gray!80!black},stringstyle=\color{black},showstringspaces=false,columns=flexible,keepspaces=true,tabsize=3,frame=single,framerule=0.4pt,rulecolor=\color{black},xleftmargin=2mm,framexleftmargin=1mm,aboveskip=2mm,belowskip=2mm,breaklines=true,
 literate={ä}{{\"a}}1 {ö}{{\"o}}1 {ü}{{\"u}}1 {Ä}{{\"A}}1 {Ö}{{\"O}}1 {Ü}{{\"U}}1 {ß}{{\ss}}1 {–}{{--}}1 {„}{{\glqq}}1 {“}{{\grqq}}1}
\lstnewenvironment{java}[1][]{\lstset{#1}}{}
% Lückencode: größerer Zeilenabstand, Lücken mit (*\luecke[3cm]*) im Code
\lstnewenvironment{javaluecke}[1][]{\lstset{basicstyle=\linespread{1.5}\small\ttfamily,escapeinside={(*}{*)},#1}}{}

\begin{document}
\abtitel{2.5}{Verwaltung gleichartiger Objekte -- Das Array / Feld}

\begin{aufgabe}[Projekt Zufallszahlen]{1}
\textbf{a)} Ergänze die Attribute \code{zahl1} bis \code{zahl3}.\quad \textbf{b)} Weise ihnen im Konstruktor Werte von 1 bis 6 zu.\\
\textbf{c)} Und bei 100 Zahlen? Skizziere, wie ein Speicher aussehen müsste, in dem du „die i-te Zahl“ ansprechen kannst:
\lsgkariert{4}{\begin{scope}[font=\scriptsize]
\foreach \v [count=\i from 0] in {5,2,6,1,4} {\draw (10mm+\i*10mm,-7mm) rectangle ++(10mm,-7mm); \node at (15mm+\i*10mm,-10.5mm) {\v}; \node at (15mm+\i*10mm,-4.5mm) {\i};}
\node at (65mm,-10.5mm) {\dots}; \draw (70mm,-7mm) rectangle ++(10mm,-7mm); \node at (75mm,-10.5mm) {3}; \node at (75mm,-4.5mm) {99};
\node[anchor=west] at (2mm,-4.5mm) {i:}; \node[anchor=west,align=left] at (85mm,-9mm) {ein Name für alle 100 Zahlen,\\Zugriff über die Nummer (Index)};
\end{scope}}
\end{aufgabe}

\begin{minipage}[t]{0.7\linewidth}
\begin{merke}[Felder]\lueckentext
Felder (engl. \lsg[3cm]{Arrays}) dienen der Speicherung einer \lsg[3cm]{festen} \lsg[3cm]{Anzahl} von Werten \lsg[4.5cm]{gleichen Datentyps}, die über einen Index (Positionszahl) adressiert werden können.

Dadurch kann z.\,B. mit einer Wiederholung auf alle Elemente zugegriffen werden, ohne jedes einzeln zu behandeln.
\end{merke}
\end{minipage}\hfill
\begin{minipage}[t]{0.27\linewidth}
\vspace{3mm}
\begin{tabular}{|c|c|}\hline
Index i & \code{textFeld[i]}\\\hline
0 & \code{"Hallo "}\\\hline
1 & \code{"Welt, "}\\\hline
2 & \code{"ein "}\\\hline
3 & \code{"Satz."}\\\hline
\end{tabular}
\end{minipage}

\begin{aufgabe}{2}
Übertrage den Code in das Projekt (TODO 2) und beschreibe die Funktion der Zeilen.
\end{aufgabe}
\begin{javaluecke}
class Zufallszahlen {
   private int[] zahlenFeld;              // (*\lsg[7.5cm]{Deklaration: Feld für int-Werte}*)
   Zufallszahlen() {
      zahlenFeld = new int[8];            // (*\lsg[7.5cm]{Erzeugung: 8 Zellen, alle mit 0 belegt}*)
      zahlenFeld[0] = 5;                  // (*\lsg[7.5cm]{erste Zelle (Index 0) bekommt 5}*)
      zahlenFeld[7] = Random.randint(1, 6); // (*\lsg[7.1cm]{letzte Zelle bekommt Zufallszahl 1--6}*)
   }
}
\end{javaluecke}

\begin{minipage}[t]{0.7\linewidth}
\begin{merke}\lueckentext
\begin{itemize}
\item Felder beginnen ihren Index immer bei \lsg[1.5cm]{0}.
\item Die Größe eines Feldes liefert \code{<Feldname>.length}. Hier also: \lsg[4cm]{\code{zahlenFeld.length}}. Das Array ist \lsg[1.5cm]{8} Zellen groß.
\item Was passiert bei \code{zahlenFeld[8]}?\lsgrest{Laufzeitfehler (Index zu groß)}
\end{itemize}
\end{merke}
\end{minipage}\hfill
\begin{minipage}[t]{0.27\linewidth}
\vspace{3mm}
\begin{tabular}{|c|c|}\hline
Index i & \code{zahlenFeld[i]}\\\hline
0 & 5\\\hline
1 & 0\\\hline
\rule[-1.2mm]{0pt}{5.2mm}2 & \lsgtext{0}\\\hline
\rule[-1.2mm]{0pt}{5.2mm}\lsgtext{3} & \lsgtext{0}\\\hline
\rule[-1.2mm]{0pt}{5.2mm}\lsgtext{4} & \lsgtext{0}\\\hline
\rule[-1.2mm]{0pt}{5.2mm}\lsgtext{5} & \lsgtext{0}\\\hline
\rule[-1.2mm]{0pt}{5.2mm}\lsgtext{6} & \lsgtext{0}\\\hline
\rule[-1.2mm]{0pt}{5.2mm}\lsgtext{7} & \lsgtext{1 bis 6}\\\hline
\end{tabular}
\end{minipage}

\newpage
Fülle die Tabelle auf Seite 1. Warum ist \code{zahlenFeld[7]} die letzte Zelle?
\lsglinien{2}{Der Index beginnt bei 0. Bei 8 Zellen laufen die Indizes von 0 bis 7, der letzte ist also 8 $-$ 1 = 7.}
{\lueckentext Letzte Zelle bei beliebiger Feldgröße:\lsgrest{\code{zahlenFeld[zahlenFeld.length - 1]}}}

\begin{aufgabe}{3}
\textbf{a)} Weise allen Zellen zufällige Werte zu (Wiederholung mit fester Anzahl).\quad
\textbf{b)} 100 Zellen, jede mit der Summe zweier Würfel.\\
Was änderst du an der Schleife aus a), damit sie bei jeder Feldgröße stimmt?
\lsglinien{1}{Bedingung \code{i < zahlenFeld.length} statt einer festen Zahl.}
\end{aufgabe}

\textbf{1:n-Beziehung als Sammlung umsetzen -- Arrays mit Referenzen}

{\centering
\begin{tikzpicture}[x=1cm,y=1cm]
\node[draw,fill=cyan!15,minimum height=6mm] (m) at (0,0) {Mannschaft};
\node[draw,fill=cyan!15,minimum height=6mm] (s) at (5,0) {Spieler};
\draw (m) -- (s) node[midway,above] {\small $<$ spielt in} node[pos=0.08,below] {\small 1} node[pos=0.92,below] {\small n};
\foreach \i/\n in {0/paul,1/tim,2/maja,3/tina} {
  \node[draw,fill=red!10,minimum width=13mm,minimum height=5mm] (z\i) at (6.6+\i*2.2,0.3) {\small\bfseries \i};
  \node[draw,fill=red!10,font=\scriptsize\ttfamily,inner sep=2pt] (o\i) at (6.6+\i*2.2,-1.1) {\n: Spieler};
  \draw[-{Stealth}] (z\i.south) -- (o\i.north);
}
\end{tikzpicture}
\par}
{\lueckentext Auch Referenzen können in Arrays verwaltet werden. Dabei verweist jede Zelle auf ein \lsg[3.5cm]{Objekt} der anderen \lsg[3.5cm]{Klasse}, hier: \lsg[3.5cm]{\code{Spieler}}.\par}

\begin{hinweis}
\lueckentext\textbf{Beachte:} Alle Elemente eines Feldes müssen denselben \lsg[3.5cm]{Datentyp} besitzen, dieser kann aber beliebig gewählt werden: auch \code{boolean}, \code{String}, \code{Rectangle}, \ldots
\end{hinweis}

\begin{merke}[Syntaxübersicht]\small
\begin{tabular}{@{}l l@{}}
\code{<Sichtbarkeit> <Datentyp>[] <Feldname>;} & Deklaration: \code{[]} nach dem Datentyp macht ein Feld daraus.\\
\code{<Feldname> = new <Datentyp>[<Groesse>];} & Erzeugung mit fester Größe (nicht veränderbar).\\
\code{<Feldname>[<Index>]} & Zugriff auf die Zelle an dieser Position.\\
\end{tabular}

\textbf{Beispiele:}\\
lesen: \code{int eineZahl = zahlenFeld[0];}\hfill schreiben: \code{zahlenFeld[0] = eineZahl;}\\
Methode aufrufen: \code{kreisFeld[0].tuEtwas();}\hfill Attribut: \code{kreisFeld[0].xPosition = 100;}
\end{merke}

\begin{minipage}[t]{0.74\linewidth}
\begin{merke}[Mehrdimensionale Arrays]\lueckentext
Arrays können mehrere Dimensionen haben und so eine \lsg[3.5cm]{Tabelle (Raster)} aufspannen.\\
Deklaration und Erzeugung:\lsgrest{\code{Rectangle[][] fuellung}}
\lsgrest{\code{fuellung = new Rectangle[5][5];}}
Zugriff auf Spalte x, Zeile y:\lsgrest{\code{fuellung[x][y]}}
\end{merke}
\end{minipage}\hfill
\begin{minipage}[t]{0.23\linewidth}
\vspace{1mm}\centering
\begin{tikzpicture}[x=6mm,y=6mm,font=\scriptsize]
\ifloesung\fill[red!25] \foreach \x/\y in {1/1,3/1,1/3,2/4,3/3} {(\x,-\y) rectangle (\x+1,-\y-1)};\fi
\draw[step=6mm] (0,0) grid (5,-5);
\foreach \i in {0,...,4} {\node at (\i+0.5,0.4) {\i}; \node at (-0.4,-\i-0.5) {\i};}
\node at (2.5,1.05) {\code{x} $\rightarrow$}; \node[rotate=-90] at (-1.05,-2.5) {\code{y} $\rightarrow$};
\end{tikzpicture}
\end{minipage}

\begin{aufgabe}[Projekt Mosaik]{4}
\textbf{a)} TODO 1 und 2 in \code{Mosaik1}: zweidimensionales Array passender Größe erzeugen und in \code{neuesMosaik()} füllen.\\
\textbf{b)} Auf \code{Mosaik2} wechseln, TODO 3--5 (Smiley). Markiere die Smiley-Zellen \code{fuellung[x][y]} im Raster oben.\lsgtext{\ (z.\,B. [1][1], [3][1], [1][3], [2][4], [3][3])}
\end{aufgabe}
\end{document}
;
%            füllt alle Lücken rot aus; siehe \lsg, \lsgrest, \lsglinien,
%            \lsgfeld, \lsgkariert, \lsgkarobox, \lsgtext, \lsgoder)
%  Lösungseinträge belegen nie mehr Platz als die Lücke im Blatt.
% =====================================================================
\usepackage[a4paper,left=18mm,right=18mm,top=26mm,bottom=17mm,headheight=15mm,headsep=4mm,footskip=8mm]{geometry}
\usepackage{iftex}
\ifPDFTeX
  \usepackage[T1]{fontenc}
  \usepackage[utf8]{inputenc}
  \usepackage{helvet}
  \usepackage[scaled=0.86]{beramono}
  \renewcommand{\familydefault}{\sfdefault}
\else
  \usepackage{fontspec}
  \setmainfont{texgyreheros}[Extension=.otf,UprightFont=*-regular,BoldFont=*-bold,ItalicFont=*-italic,BoldItalicFont=*-bolditalic]
  \setmonofont{DejaVuSansMono}[Extension=.ttf,UprightFont=*,BoldFont=*-Bold,ItalicFont=*-Oblique,BoldItalicFont=*-BoldOblique,Scale=0.86]
\fi
\usepackage[ngerman,shorthands=off]{babel}
\usepackage{xcolor,graphicx,tabularx,array,enumitem,fancyhdr,tikz,listings,multicol,calc,etoolbox,ragged2e,colortbl,xspace,pifont}
\usepackage[most]{tcolorbox}
\usetikzlibrary{arrows.meta,positioning,calc,decorations.pathreplacing}
\usepackage[hidelinks]{hyperref}
\setlength{\parindent}{0pt}
\setlength{\parskip}{3pt}
\setlist{nosep,leftmargin=*}

% ---------- Lösungsschalter ----------
\newif\ifloesung
\ifdefined\mitloesung\loesungtrue\fi
\newcommand{\lsgfarbe}{\color{red!75!black}}

% ---------- Kopf- und Fußzeile (einheitliche Form) ----------
\newcommand{\lehrkraft}{Hau}
\newcommand{\themenbereich}{2. Objektorientierte Programmierung}
\newcommand{\abnummer}{}
\pagestyle{fancy}\fancyhf{}
\renewcommand{\headrulewidth}{0.4pt}
\fancyhead[L]{\small Klasse 10\\Datum:}
\fancyhead[C]{\small Themenbereich:\\\themenbereich}
\fancyhead[R]{\small \lehrkraft\ifloesung\\{\lsgfarbe\bfseries Lösung}\fi}
\fancyfoot[C]{\scriptsize edu-mrh.de\,·\,Informatik 10\,·\,Blatt \abnummer\,·\,Seite \thepage}
\newcommand{\abtitel}[2]{\renewcommand{\abnummer}{#1}\begin{center}\Large\bfseries #1\quad\underline{#2}\end{center}\vspace{-1mm}}

% ---------- Boxen ----------
\newenvironment{aufgabe}[2][]{\begin{tcolorbox}[enhanced,breakable,colback=white,colframe=black,boxrule=0.7pt,arc=3mm,left=2mm,right=2mm,top=1.5mm,bottom=1.5mm,before skip=2.5mm,after skip=2.5mm]\textbf{Aufgabe #2}\ifstrempty{#1}{}{ \textit{#1}}\par\vspace{0.6mm}}{\end{tcolorbox}}
\newtcolorbox{merke}[1][Merke]{enhanced,breakable,colback=green!5,colframe=green!40!black,boxrule=0.8pt,arc=2mm,left=2mm,right=2mm,top=1.5mm,bottom=1.5mm,before skip=2.5mm,after skip=2.5mm,title={#1},fonttitle=\bfseries,coltitle=white}
\newtcolorbox{hinweis}{enhanced,breakable,colback=yellow!12,colframe=orange!70!black,boxrule=0.6pt,arc=2mm,left=2mm,right=2mm,top=1mm,bottom=1mm,before skip=2mm,after skip=2mm}
\newcommand{\web}[1]{{\small\color{gray}\textit{Website: #1}}}

% ---------- Lücken, Linien, Karos ----------
% Einheitliche Maße für alle Blätter der Sequenz:
%  \linien{n}       n Schreiblinien, Abstand 8 mm, erste Linie 8 mm unter dem Text
%  \kariert{n}      Karoraster 5 mm, n Kästchen hoch, nur ganze Kästchen, mittig
%  \karobox{b}{h}   Karoraster b x h, auf ganze Kästchen gerundet
%  \luecke[l]       Lücke im Fließtext (Standard 4,5 cm, Fachbegriffe mind. 3,5 cm)
%  (jeweils mit Lösungsvariante \lsg…, siehe Kopf der Datei)
%  \restzeile       Lücke bis zum Zeilenende
%  \lueckentext     am Anfang eines Absatzes/Kastens mit Lücken: größerer
%                   Zeilenabstand und Flattersatz, damit man in die Lücken schreiben kann
%  \linienfeld{n}   umrahmtes Feld mit n Schreiblinien (z. B. für Java-Code)
%  \schreibzeile    Strut für Tabellenzellen, in die geschrieben wird (8 mm hoch)
\newlength{\zeilenabstand}\setlength{\zeilenabstand}{8mm}
\newlength{\karo}\setlength{\karo}{5mm}
\tikzset{karo/.style={draw=gray!50,line width=0.3pt},schreiblinie/.style={draw=black!65,line width=0.35pt}}
\newlength{\lsgbreite}
% Lücke: \luecke[Breite] / mit Lösung \lsg[Breite]{Text} (zu langer Text wird verkleinert)
\newcommand{\lsg}[2][4.5cm]{\leavevmode\ifloesung\settowidth{\lsgbreite}{\,#2\,}%
  \ifdim\lsgbreite>#1\relax\rlap{\raisebox{0.8pt}{\resizebox{#1}{!}{\lsgfarbe\,#2\,}}}%
  \else\rlap{\raisebox{0.8pt}{\lsgfarbe\,#2}}\fi\fi\rule[-0.6ex]{#1}{0.4pt}\xspace}
\newcommand{\luecke}[1][4.5cm]{\lsg[#1]{}}
% Lücke bis Zeilenende: \restzeile / \lsgrest{Text}
\newcommand{\lsgrest}[1]{\leavevmode\unskip\hspace{1.5mm}\ifloesung\rlap{\raisebox{0.8pt}{\lsgfarbe\,#1}}\fi\leaders\hbox{\rule[-0.6ex]{0.5pt}{0.4pt}}\hfill\null\par}
\newcommand{\restzeile}{\lsgrest{}}
% nur in der Lösung sichtbar (Platz bleibt frei): \lsgtext{…}; Alternativen: \lsgoder{Lösung}{Blatt}
\newcommand{\lsgtext}[1]{\ifloesung\leavevmode{\lsgfarbe#1}\fi}
\newcommand{\lsgoder}[2]{\ifloesung\leavevmode#1\else#2\fi}
\newcommand{\lueckentext}{\linespread{1.5}\selectfont\raggedright}
\newcommand{\schreibzeile}{\rule[-2.8mm]{0pt}{8mm}}
% Schritt-Nummern wie im Code und auf der Website: \nr{1} = ①
\newcommand{\nr}[1]{\ding{\numexpr171+#1\relax}}
% Schreiblinien: \linien{n} / \lsglinien{n}{Text} (Text liegt auf den Linien)
\newcommand{\lsglinien}[2]{\par\vspace{0.5mm}\noindent\begin{tikzpicture}
  \path (0,0) (\linewidth-0.4pt,-#1\zeilenabstand-1.5mm);
  \foreach \i in {1,...,#1}{\draw[schreiblinie] (0.2pt,-\i\zeilenabstand) -- (\linewidth-0.2pt,-\i\zeilenabstand);}
  \ifloesung\node[overlay,anchor=north west,inner sep=0pt] at (1mm,-\zeilenabstand+1pt+5mm) {\parbox[t]{\linewidth-2mm}{\lsgfarbe\raggedright\setlength{\baselineskip}{\zeilenabstand}\rule{0pt}{5mm}#2}};\fi
  \end{tikzpicture}\par}
\newcommand{\linien}[1]{\lsglinien{#1}{}}
% Codefeld mit Linien: \linienfeld{n} / \lsgfeld{n}{Code, Zeilen mit \\ getrennt}
\newcommand{\lsgfeld}[2]{\noindent\begin{tikzpicture}
  \draw[line width=0.4pt] (0.2pt,0) rectangle (\linewidth-0.2pt,-#1\zeilenabstand-3mm);
  \foreach \i in {1,...,#1}{\draw[schreiblinie] (2mm,-\i\zeilenabstand) -- (\linewidth-2mm,-\i\zeilenabstand);}
  \ifloesung\node[overlay,anchor=north west,inner sep=0pt] at (2.5mm,-\zeilenabstand+1pt+5mm) {\parbox[t]{\linewidth-5mm}{\lsgfarbe\ttfamily\small\setlength{\baselineskip}{\zeilenabstand}\rule{0pt}{5mm}#2}};\fi
  \end{tikzpicture}}
\newcommand{\linienfeld}[1]{\lsgfeld{#1}{}}
% Karoraster; Lösungszeichnung als TikZ-Code, Ursprung oben links, y nach unten negativ:
% \kariert{n} / \lsgkariert{n}{TikZ},  \karobox{b}{h} / \lsgkarobox{b}{h}{TikZ}
\newcommand{\karoraster}[3]{\begin{tikzpicture}\draw[karo] (0,0) grid[step=\karo] (#1\karo,#2\karo);
  \ifloesung\begin{scope}[overlay,shift={(0,#2\karo)},red!75!black,text=red!75!black,line width=0.8pt]#3\end{scope}\fi\end{tikzpicture}}
\newcommand{\lsgkariert}[2]{\par\vspace{1.5mm}\noindent\pgfmathtruncatemacro{\karoX}{floor(\linewidth/\karo)}%
  \makebox[\linewidth]{\karoraster{\karoX}{#1}{#2}}\par\vspace{1mm}}
\newcommand{\kariert}[1]{\lsgkariert{#1}{}}
\newcommand{\lsgkarobox}[3]{\pgfmathtruncatemacro{\karoX}{round((#1)/\karo)}\pgfmathtruncatemacro{\karoY}{round((#2)/\karo)}\karoraster{\karoX}{\karoY}{#3}}
\newcommand{\karobox}[2]{\lsgkarobox{#1}{#2}{}}
% Lösungs-Klassenkarte (für Zeichnungen im Karo): \lsgkarte{Name}{Attribute}{Methoden}
\newcommand{\lsgkarte}[3]{{\lsgfarbe\setlength{\arrayrulewidth}{0.6pt}\arrayrulecolor{red!75!black}\begin{tabular}{|l|}\hline\textbf{#1}\\\hline #2\\\hline #3\\\hline\end{tabular}}}
% Lösungs-Objektkarte (in TikZ als Knoteninhalt): \node[objekt] {\lsgobjekt{name : Klasse}{attr = wert\\…}};
\newcommand{\lsgobjekt}[2]{{\lsgfarbe\small\arrayrulecolor{red!75!black}\begin{tabular}{@{}l@{}}\textbf{#1}\\\hline #2\end{tabular}}}
\tikzset{objekt/.style={draw,rounded corners=3mm,inner sep=2mm,anchor=north west}}
\newcommand{\klassenkarte}[3]{\begin{tabular}{|l|}\hline\textbf{#1}\\\hline #2\\\hline #3\\\hline\end{tabular}}
\newcommand{\code}[1]{\texttt{#1}}

% ---------- Java-Listings ----------
\lstset{language=Java,basicstyle=\ttfamily\small,keywordstyle=\bfseries,commentstyle=\color{gray!80!black},stringstyle=\color{black},showstringspaces=false,columns=flexible,keepspaces=true,tabsize=3,frame=single,framerule=0.4pt,rulecolor=\color{black},xleftmargin=2mm,framexleftmargin=1mm,aboveskip=2mm,belowskip=2mm,breaklines=true,
 literate={ä}{{\"a}}1 {ö}{{\"o}}1 {ü}{{\"u}}1 {Ä}{{\"A}}1 {Ö}{{\"O}}1 {Ü}{{\"U}}1 {ß}{{\ss}}1 {–}{{--}}1 {„}{{\glqq}}1 {“}{{\grqq}}1}
\lstnewenvironment{java}[1][]{\lstset{#1}}{}
% Lückencode: größerer Zeilenabstand, Lücken mit (*\luecke[3cm]*) im Code
\lstnewenvironment{javaluecke}[1][]{\lstset{basicstyle=\linespread{1.5}\small\ttfamily,escapeinside={(*}{*)},#1}}{}

\begin{document}
\abtitel{2.0.2}{Vererbung: Wiederholung und Vertiefung}

\begin{aufgabe}[Das gallische Dorf]{0}
In einem gallischen Dorf leben Druiden, Barden, Hinkelsteinverkäufer und ein Häuptling. Alle Dorfbewohner haben einen Namen und Muskelkraft, die sie durch das Trinken eines Zaubertranks verzehnfachen können. Außerdem können sie Fische werfen. Nur Druiden können einen Zaubertrank brauen, Barden können singen und nur ein Häuptling kann seine Schildträger rufen. Hinkelsteinverkäufer haben hingegen ein Vermögen. Außerdem gibt es verschiedene Tiere, die Laute von sich geben können. Hunde haben zusätzlich einen Namen, während Wildschweine weglaufen können.\\
\textbf{a)} Zeichne ein Klassendiagramm, das diese Situation modelliert.\quad \textbf{b)} Programmieren: Website.
\end{aufgabe}
\lsgkariert{11}{\begin{scope}[font=\scriptsize,every node/.style={anchor=north west,inner sep=0pt}]
\node (db) at (35mm,-3mm) {\lsgkarte{Dorfbewohner}{name\\muskelkraft}{zaubertrankTrinken()\\fischeWerfen()}};
\node (dr) at (3mm,-31mm) {\lsgkarte{Druide}{\rule{0pt}{2mm}}{zaubertrankBrauen()}};
\node (ba) at (33mm,-31mm) {\lsgkarte{Barde}{\rule{0pt}{2mm}}{singen()}};
\node (hp) at (51mm,-31mm) {\lsgkarte{Häuptling}{\rule{0pt}{2mm}}{schildtraegerRufen()}};
\node (hv) at (82mm,-31mm) {\lsgkarte{Hinkelsteinverkäufer}{vermoegen}{\rule{0pt}{2mm}}};
\node (ti) at (132mm,-3mm) {\lsgkarte{Tier}{\rule{0pt}{2mm}}{lautGeben()}};
\node (hu) at (118mm,-31mm) {\lsgkarte{Hund}{name}{\rule{0pt}{2mm}}};
\node (ws) at (140mm,-31mm) {\lsgkarte{Wildschwein}{\rule{0pt}{2mm}}{weglaufen()}};
\end{scope}
\begin{scope}[-{Triangle[open,length=2.5mm]}]
\foreach \k in {dr,ba,hp,hv} \draw (\k.north) -- (db.south);
\foreach \k in {hu,ws} \draw (\k.north) -- (ti.south);
\end{scope}}

\vspace{1mm}
Fehlermeldung beim ersten \code{new Barde(...)} und ihre Lösung:
\lsglinien{2}{Barde hat keinen passenden Konstruktor (Konstruktoren werden nicht vererbt). Lösung: eigener Konstruktor \code{Barde(String nameNeu, int muskelkraftNeu) \{ super(nameNeu, muskelkraftNeu); \}}}

\begin{aufgabe}{1}
\RaggedRight Zeichne ein vereinfachtes Klassendiagramm (nur Klassennamen) zu: \code{Tier}, \code{Wirbeltiere}, \code{WirbelloseTiere}, \code{Fische}, \code{Amphibien}, \code{Saeugetiere}, \code{Voegel}, \code{Reptilien}, \code{Insekten}, \code{Spinnen}.
\end{aufgabe}
\lsgkariert{5}{\begin{scope}[font=\scriptsize,every node/.style={draw,inner sep=1.5pt}]
\node (t) at (60mm,-3.5mm) {Tier};
\node (w) at (40mm,-12.5mm) {Wirbeltiere}; \node (l) at (120mm,-12.5mm) {WirbelloseTiere};
\foreach \n/\x in {Fische/8,Amphibien/30,Saeugetiere/54,Voegel/76,Reptilien/96} \node (\n) at (\x mm,-21.5mm) {\n};
\node (i) at (112mm,-21.5mm) {Insekten}; \node (sp) at (134mm,-21.5mm) {Spinnen};
\end{scope}
\begin{scope}[-{Triangle[open,length=1.8mm]},line width=0.6pt]
\draw (w) -- (t); \draw (l) -- (t);
\foreach \n in {Fische,Amphibien,Saeugetiere,Voegel,Reptilien} \draw (\n) -- (w);
\draw (i) -- (l); \draw (sp) -- (l);
\end{scope}}

\begin{merke}[Vererbungshierarchie, Ober- und Unterklassen]\lueckentext
Bei der Vererbung gibt es Unter- und Oberklassen. Eine \lsg{Unterklasse} besitzt automatisch alle Attribute und Methoden der \lsg{Oberklasse}. Es werden zwei Richtungen unterschieden: Wird eine Unterklasse erstellt, so spricht man von einer \lsg{Spezialisierung} der Oberklasse. Umgekehrt ist die Oberklasse eine \lsg{Generalisierung} der Unterklasse.

Ist eine Unterklasse gleichzeitig Oberklasse einer weiteren Klasse, spricht man von \textbf{mehrstufiger Vererbung}; es entsteht eine \lsg[5cm]{Vererbungshierarchie}.
\end{merke}

\newpage
\begin{merke}[Vererbung und Überschreiben]\lueckentext
Konstruktoren werden \lsg[3.5cm]{nicht} vererbt. Die Unterklasse ruft den Konstruktor der Oberklasse mit \lsg[4cm]{\code{super(…)}} auf.

Definiert man eine Methode in der Unterklasse erneut (gleicher Name, gleiche Parameter, gleicher Rückgabetyp), so wird die ursprüngliche Methode \lsg{überschrieben}. Sie hat in der Unterklasse eine \textbf{andere} Funktionsweise.

Die Annotation \code{@Override} vor der Methode sagt dem Compiler: Diese Methode überschreibt eine gleichnamige Methode der Oberklasse. Gibt es diese nicht (z.\,B. Tippfehler), meldet der Compiler einen \lsg[3.5cm]{Fehler}. Ohne \code{@Override} entsteht stattdessen stillschweigend \lsg[6cm]{eine neue, zusätzliche Methode}.
\end{merke}

\begin{aufgabe}[Projekt Aquarium]{2}
\textbf{a)} Nenne die Attribute (mit Datentyp), die die Klasse \code{Fisch} bereitstellt.
\lsglinien{2}{\code{String art}, \code{boolean tarnung} (Startwert \code{false}), \code{int geschwindigkeit}}
\textbf{b)} Beschreibe den Konstruktor von \code{Fisch}: Parameter und Wirkung.
\lsglinien{2}{\code{Fisch(String art, int x, int y)}: setzt Position und Bild mit \code{super(x, y, …)}, speichert die Art, wählt eine zufällige Geschwindigkeit 1--4 und vergrößert das Bild.}
\textbf{c)}--\textbf{e)} auf der Website. Attribute deiner Klasse \code{Hai} mit Datentyp und Startwert:
\lsglinien{1}{\code{int alter = 1;}\quad \code{boolean hungrig = true;}}
\end{aufgabe}

\begin{aufgabe}[Fehlersuche]{3}
Im folgenden Code haben sich Fehler eingeschlichen. Markiere sie im Code, gib in der Tabelle die Art des Fehlers an (Syntax, Datentyp, Logik) und schreibe die Korrektur daneben.
\end{aufgabe}
\begin{minipage}[t]{0.5\linewidth}
\begin{java}[numbers=left,numberstyle=\scriptsize\color{gray!70!black},numbersep=4pt,xleftmargin=5mm]
class Katze
{
   String name;
   String alter;
   boolean geschlecht;

   public Katze(String nameN, int alterN)
   {
      name = "Minka";
      alter = alterN
      geschlecht = "weiblich";
   }
}
\end{java}
\end{minipage}\hfill
\begin{minipage}[t]{0.46\linewidth}
\vspace{0pt}\renewcommand{\arraystretch}{1}
\begin{tabularx}{\linewidth}{|c|p{1.7cm}|X|}\hline
\footnotesize Zeile & \footnotesize Art & \footnotesize Korrektur\\\hline
\schreibzeile\lsgtext{4} & \lsgtext{Datentyp} & \lsgtext{\code{int alter;}}\\\hline
\schreibzeile\lsgtext{9} & \lsgtext{Logik} & \lsgtext{\code{name = nameN;}}\\\hline
\schreibzeile\lsgtext{10} & \lsgtext{Syntax} & \lsgtext{\code{;} am Zeilenende}\\\hline
\schreibzeile\lsgtext{11} & \lsgtext{Datentyp} & \lsgtext{\code{geschlecht = false;}}\\\hline
\end{tabularx}
\end{minipage}
\end{document}
